\documentclass[11pt]{article}

\usepackage[utf8]{inputenc}
\usepackage[T1]{fontenc}
\usepackage{lmodern}
\usepackage{geometry}
\usepackage{amsmath,amssymb,amsfonts,amsthm,mathtools}
\usepackage{bm}
\usepackage{enumitem}
\usepackage{microtype}
\usepackage{hyperref}
\usepackage{titlesec}
\usepackage{booktabs}
\usepackage{array}
\usepackage{setspace}
\usepackage{mathrsfs}
\usepackage{physics}

\geometry{margin=1in}
\hypersetup{colorlinks=true, linkcolor=black, urlcolor=blue, citecolor=black}
\setlist[enumerate]{topsep=0.4em,itemsep=0.35em}
\setlist[itemize]{topsep=0.3em,itemsep=0.25em}
\titleformat{\section}{\large\bfseries}{\thesection.}{0.5em}{}
\titleformat{\subsection}{\normalsize\bfseries}{\thesubsection.}{0.5em}{}
\setstretch{1.06}

\newcommand{\Aether}{\AE{}ther}
\newcommand{\Aflow}{\AE{}ther-flow}
\newcommand{\TheoryName}{\Aflow{} Interpretation of Relativity}
\newcommand{\TheoryProgram}{\Aether{} / \Aflow{} framework}

\newtheorem{definition}{Definition}
\newtheorem{proposition}{Proposition}
\newtheorem{remark}{Remark}
\newtheorem{corollary}{Corollary}

\title{\textbf{The \TheoryName{}}\\[0.4em]
\large Flagship Public Article}
\author{}
\date{}

\begin{document}

\maketitle

\begin{abstract}
This article is the single-entry public introduction to \TheoryName{}. Its job is not to replace the active exact-closure manuscript sequence, but to make that sequence readable to first-time visitors, technically trained readers, and future collaborators. The central public message is simple. The flagship deliverable of the repository is already the overview-first exact-closure package centered on the \emph{Exact-Closure Sequence Overview}. Within that package, the effective gravitational dynamics are adopted to be exactly those of general relativity with universal matter coupling, while the \Aether{} / \Aflow{} ontology supplies the interpretive and conceptual setting.

The article makes five points explicit. First, the active theory statement is exact relativistic closure, not a low-energy deformation of GR. Second, the \Aether{} / \Aflow{} language is retained in disciplined ontological form through the notions of underlying substrate, intrinsic ordered motion, observed three-dimensional space as local experiential slice, and S-time as experienced order of change. Third, the public claim boundary is strict: GR is adopted exactly in the active package, but a completed first-principles substrate derivation is not yet claimed. Fourth, the six core manuscripts should be read as one coordinated package after the overview. Fifth, the larger derivational bridge archive remains available for transparency and future work, but it is downstream and optional relative to the flagship package.
\end{abstract}

\section{Introduction}

The repository now has a stable public-facing theory package. That fact changes the right mode of presentation. The project no longer needs to be introduced primarily as ontology awaiting its first disciplined theory statement. The exact theory statement already exists in the overview-first exact-closure sequence. What is now needed is a public article that explains that package cleanly, states the claim boundary honestly, and helps readers know where to start.

This article therefore has a narrower role than the companion manuscripts. The \emph{Exact-Closure Sequence Overview} remains the canonical front door to the package itself, and the \emph{Exact-Closure Note} remains the shortest standalone anchor. The present article is instead a public-facing synthesis. It explains what the active package is, what it does and does not claim, how the companion papers fit together, and why the larger derivational archive should be read only afterward.

\paragraph{Framework and claim boundary.}
\TheoryProgram{} denotes the broader \Aether{} / \Aflow{} framework built on the claims that the \Aether{} is the underlying four-dimensional substrate of reality, the \Aflow{} is its intrinsic ordered motion, observed three-dimensional space is the local experiential slice of that deeper substrate, S-time is the experienced order of change arising from matter, light, and the \Aflow{}, and observed expansion is the three-dimensional appearance of deeper four-dimensional ordered motion. Within that framework, \TheoryName{} denotes the exact-closure theory in which the effective gravitational dynamics are adopted to be exactly Einsteinian with universal matter coupling. In the active sequence, \emph{adoption} means use of that established relativistic dynamics without claiming substrate derivation, while \emph{derivation} is reserved for a first-principles recovery from explicit substrate variables.

\section{Public Role of This Article}

\begin{definition}[Flagship public article]
A flagship public article is a deployment-facing manuscript that introduces the active exact-closure package in one readable document, preserves the fixed claim boundary of the package, and directs readers into the full manuscript sequence without reopening the underlying theory line.
\end{definition}

\begin{proposition}[Sequence-preserving role]
The present article is secondary to the exact-closure sequence overview in authority, but primary as a public orientation document for readers who need a single coherent entry article.
\end{proposition}

This distinction matters. If the public article were allowed to replace the exact-closure overview, the package could drift into a looser or more promotional presentation. If it were omitted entirely, first-time readers would be left to infer the status of the project from a distributed manuscript tree. The correct middle position is to let the package remain the source of truth while the present article functions as its public guide.

\section{Ontology in Disciplined Form}

The conceptual core of \TheoryProgram{} is ontological rather than deformation-first. The active vocabulary is:
\begin{itemize}
    \item \Aether{}: the underlying four-dimensional substrate of reality;
    \item \Aflow{}: the intrinsic ordered motion of that substrate;
    \item observed three-dimensional space: the local experiential slice of the deeper substrate;
    \item S-time: the experienced order of change arising from matter, light, and the \Aflow{};
    \item observed expansion: the three-dimensional appearance of deeper four-dimensional ordered motion.
\end{itemize}

These terms are intended to discipline the interpretive proposal, not to overstate the dynamical derivation. The \Aflow{} is not introduced as a naive vector wind in ordinary observed space. The observed world is not identified with the whole ontology. The point of the framework is instead that the observer-accessible relativistic world may be interpreted as the experienced appearance of deeper ordered substrate structure.

\begin{remark}
At the active public stage, this ontology is real as interpretive structure, but it is not yet advertised as a completed microphysical derivation of the relativistic sector.
\end{remark}

\section{Exact Relativistic Theory Statement}

\begin{definition}[\TheoryName{}]
\TheoryName{} is the exact-closure, GR-consistent theory of \TheoryProgram{} in which the \Aether{} / \Aflow{} ontology is retained while the effective gravitational dynamics are exactly those of general relativity with universal matter coupling.
\end{definition}

Its effective action is
\begin{equation}
S_{\mathrm{eff}}
=
\frac{c^3}{16\pi G}\int d^4x\,\sqrt{-g}\,R
+
S_{\mathrm{matter}}[g,\psi],
\label{eq:Seff}
\end{equation}
with field equations
\begin{equation}
G_{\mu\nu}
=
\frac{8\pi G}{c^4}T_{\mu\nu},
\label{eq:Einstein}
\end{equation}
the standard null condition
\begin{equation}
0=g_{\mu\nu}\,dx^\mu dx^\nu,
\label{eq:null}
\end{equation}
and the standard proper-time rule
\begin{equation}
d\tau^2=-\frac{1}{c^2}g_{\mu\nu}\,dx^\mu dx^\nu.
\label{eq:tau}
\end{equation}

\begin{proposition}[Operational identity with GR]
For fixed matter sector, initial data, and boundary conditions, the effective observable content of \TheoryName{} is exactly that of general relativity.
\end{proposition}

This proposition is the center of the exact-closure package. The project does not currently claim a verified low-energy deformation of GR. It claims that the ontology can be retained while the effective gravitational law remains exactly Einsteinian.

\section{Adoption, Interpretation, and Open Derivation}

\begin{definition}[Adoption]
Adoption is the use of an already established effective dynamics, here GR with universal matter coupling, without claiming that the dynamics have already been derived from deeper substrate variables.
\end{definition}

\begin{definition}[Derivation]
Derivation is a first-principles recovery of the effective relativistic sector from explicit substrate variables, together with whatever mode control, uniqueness, and closure arguments are needed to justify that recovery.
\end{definition}

The active public claim of \TheoryName{} is adoption with exact closure, not completed derivation. That boundary should be stated positively rather than defensively. It means that the effective theory statement is already disciplined and complete at the operational level, while the deeper foundational recovery problem remains open.

The package therefore establishes the following.
\begin{enumerate}
    \item The ontological vocabulary is fixed in disciplined form.
    \item The exact effective relativistic law is fixed as Einsteinian with universal matter coupling.
    \item The package carries no separate established low-energy non-GR observable sector.
    \item The project distinguishes clearly between interpretation, adoption, and derivation.
\end{enumerate}

The package also does \emph{not} claim the following.
\begin{enumerate}
    \item It does not claim that GR has already been derived from substrate microphysics.
    \item It does not claim that ontology by itself settles the microphysical problem.
    \item It does not claim that the larger bridge archive has already produced a completed Einstein derivation.
    \item It does not claim that the present public package must be read through the bridge archive before it can count as a theory statement.
\end{enumerate}

\section{The Flagship Exact-Closure Package}

The active public package should be read in the following fixed order.

\begin{center}
\begin{tabular}{p{0.07\textwidth} p{0.31\textwidth} p{0.52\textwidth}}
\toprule
\textbf{Order} & \textbf{Manuscript} & \textbf{Role} \\
\midrule
1 & Exact-Closure Sequence Overview & Canonical front door fixing the package role, reading order, and claim boundary. \\
2 & Exact-Closure Note & Short standalone anchor stating the exact-closure theory in its briefest complete form. \\
3 & Foundations & Ontological vocabulary, framework intent, and the exact-closure relation to GR. \\
4 & Dynamics & Effective action, matter coupling, weak-field structure, and benchmark exact law. \\
5 & Consistency & Gauge structure, degree-of-freedom counting, and effective health conditions. \\
6 & Relativistic Recovery & Exact relation to GR and local relation to SR. \\
7 & Flow Geometry & Congruence-based geometric dictionary for reading \Aflow{} inside exact relativistic structure. \\
\bottomrule
\end{tabular}
\end{center}

\begin{corollary}[Public packaging rule]
The overview-first exact-closure sequence is the flagship deliverable of the repository and should be presented before any downstream derivational continuation.
\end{corollary}

This order is not arbitrary. It places the finished effective theory statement before the open foundational program. Public accessibility and scientific honesty both improve when that order is held fixed.

\section{Relation to the Downstream Derivational Bridge Archive}

The repository also contains a large derivational bridge archive. That archive remains important. It documents genuine attempts to recover or constrain the deeper substrate route, records negative and partial results, and preserves the history of the broader research program. Yet its role in public presentation is now different from what it would have been at an earlier stage.

The bridge archive is no longer the condition for saying that the repository has a theory statement. The theory statement is already given by \TheoryName{} as the exact-closure package. The archive is therefore downstream and optional. It should be read as foundational continuation answerable to the exact-closure benchmark, not as the front door to the project.

\begin{remark}
If a future bridge route succeeds, it must recover the same exact relativistic benchmark already fixed by the flagship package. If a future bridge route fails, the project contracts back to the exact-closure theory rather than back to ontology alone.
\end{remark}

\section{Research Provenance and Collaboration Method}

The repository should also be read honestly as AI-assisted theoretical research under sustained human direction. The original conceptual direction came from the user. GPT-5.4 assisted with drafting, mathematical exploration, branch screening, organization, and revision. The resulting package is therefore neither a purely conventional solo draft nor an autonomous AI production.

This provenance matters for transparency, but it does not change the intellectual standard to which the work should be held. The manuscripts should be judged by their assumptions, equations, internal logic, and honesty about what remains open. AI assistance does not convert speculative steps into established physics, and it does not remove the need for expert human review. Its role here has been methodological support within an explicitly documented collaboration.

\section{How the Package Should Be Reviewed}

At public-release stage, the most useful technical review focuses on four questions.
\begin{enumerate}
    \item Is the claim boundary stable and explicit throughout the package?
    \item Do the equations and prose say exactly what is established, and no more?
    \item Are the package roles of overview, anchor, and companion manuscripts clear and non-overlapping?
    \item Does the presentation distinguish properly between the flagship exact-closure package and the downstream bridge archive?
\end{enumerate}

Readers who need the shortest technically disciplined path should begin with the overview and then the exact-closure note. Readers who want the modular full statement should continue through foundations, dynamics, consistency, relativistic recovery, and flow geometry in that order. Only after that should the bridge archive be used for broader foundational context.

\section{Conclusion}

The repository now has a public-facing flagship theory package. The positive result is that \TheoryName{} can already be presented as a disciplined exact relativistic theory statement built on the \Aether{} / \Aflow{} ontology, with the effective gravitational dynamics adopted exactly as those of GR. The scope of that result is the overview-first exact-closure package and its companion manuscripts, not a completed first-principles substrate derivation.

The next honest burden is therefore twofold. First, public release materials should continue to make this exact-closure package readable, accessible, and reviewable. Second, any future foundational continuation should remain clearly subordinate to that package and should state without ambiguity that deeper derivation from substrate variables remains open work rather than an achieved public claim.

\end{document}
